\documentclass[12pt]{article}
\usepackage[utf8]{inputenc}
\usepackage[T1]{fontenc}
\usepackage{hyperref}
\title{Zeem Language Resources for K-12 Curriculum Development}
\author{Compiled by Claude Sonnet 5.5 (claude-sonnet-5-5) for Andr\'as Kornai}
\date{}
\begin{document}
\maketitle

\section{Introduction}
Zeem (ISO 639-3 zem) is a Chadic language of the Afro-Asiatic family, spoken in Toro LGA, Bauchi State, Nigeria. Ethnologue classes it as dormant, with no remaining first-language users. Wikipedia instead lists about 400 speakers (2003). Joshua Project reports about 1,700 Zeem people, who increasingly prefer Hausa or English. Nothing was found that could serve as school material.

\section{Primary Educational Resources}
No primers, textbooks or school programmes were found.

\section{Linguistic and Reference Materials}
\begin{itemize}
\item Zeem language, Wikipedia: \url{https://en.wikipedia.org/wiki/Zeem_language}
\item Ethnologue entry for Zeem (subscription-limited): \url{https://www.ethnologue.com/language/zem}
\item Endangered Languages Project entry, which uses the older code zua: \url{https://endangeredlanguages.com/lang/3094}
\end{itemize}

\section{Religious and Community Resources}
Joshua Project has a people-group profile for the Zeem of Nigeria: \url{https://joshuaproject.net/people_groups/16097}

\section{Digital and Audio Resources}
No audio, video or corpus resources were found.

\section{Curriculum Development Context}
With no known first-language speakers per Ethnologue and no known teaching materials, any curriculum would be a revitalisation project drawing on archival documentation. The searches did not identify such documentation. The related variety Dyarim and the Chadic context could be investigated further.
\end{document}